\DocumentMetadata{
  pdfversion=2.0,
  pdfstandard=ua-2,
  lang=en-US,
  tagging=on,
  tagging-setup={math/setup=mathml-SE,tabsorder=structure}
}
\documentclass[11pt,letterpaper]{article}
\usepackage[margin=0.68in,includefoot]{geometry}
\usepackage{newcomputermodern}
\usepackage{amsmath,amscd,mathtools}
\usepackage{tikz,adjustbox}
\providecommand{\square}{\mdlgwhtsquare}
\usepackage[hidelinks]{hyperref}
\hypersetup{
  pdftitle={Numerical Analysis PhD Exam, May 2020},
  pdfauthor={University of Florida Department of Mathematics},
  pdfsubject={Graduate mathematics examination},
  pdfdisplaydoctitle=true,
  bookmarksopen=true
}
\setcounter{secnumdepth}{0}
\setlength{\parindent}{0pt}
\setlength{\parskip}{0.28em}
\setlength{\emergencystretch}{3em}
\setlength{\leftmargini}{2.15em}
\setlength{\leftmarginii}{2.65em}
\setlength{\leftmarginiii}{3em}
\pagestyle{empty}
\raggedbottom
\allowdisplaybreaks
\NewTaggingSocketPlug{block/list/label}{examproblem}{%
  \tagstructbegin{tag=Lbl}\tagstructend%
  \tagstructbegin{tag=\LBody}%
  \tagstructbegin{tag=\ProblemHeadingTag,title-o={\ProblemHeadingText}}%
  \tagmcbegin{tag=\ProblemHeadingTag,actualtext={\ProblemHeadingText}}%
  #2%
  \tagmcend%
  \tagstructend%
}
\newenvironment{examproblems}[2]{%
  \def\ProblemHeadingTag{#2}%
  \AssignTaggingSocketPlug{block/list/label}{examproblem}%
  \begin{enumerate}%
  \setcounter{enumi}{#1}%
  \renewcommand{\labelenumi}{\textbf{\theenumi.}}%
  \setlength{\topsep}{0.35em}%
  \setlength{\partopsep}{0pt}%
  \setlength{\itemsep}{0.48em}%
  \setlength{\parsep}{0.18em}%
}{\end{enumerate}}
\newenvironment{examsubparts}{%
  \AssignTaggingSocketPlug{block/list/label}{default}%
  \begin{enumerate}%
  \setlength{\topsep}{0.18em}%
  \setlength{\partopsep}{0pt}%
  \setlength{\itemsep}{0.16em}%
  \setlength{\parsep}{0.1em}%
}{\end{enumerate}}
\newenvironment{examinstructions}{%
  \par\begingroup\itshape
}{%
  \par\endgroup\vspace{0.18em}
}
\makeatletter
\renewcommand\subsection{\@startsection{subsection}{2}{0pt}%
  {-1.25ex plus -.25ex minus -.1ex}%
  {0.55ex plus .1ex}%
  {\normalfont\large\bfseries}}
\renewcommand{\@maketitle}{%
  \newpage\null\vskip -1em%
  {\footnotesize\bfseries
    \href{https://www.ufl.edu/}{University of Florida}, \href{https://math.ufl.edu/}{Department of Mathematics}\par}%
  \vskip -0.5em%
  \tagstructbegin{tag=H1,title={Numerical Analysis PhD Exam, May 2020}}%
  \tagpdfparaOff%
  \tagmcbegin{tag=H1}%
    {\LARGE\bfseries\href{https://gma.math.ufl.edu/exams/numerical-analysis-phd/}{Numerical Analysis PhD Exam}, May 2020\par}%
  \tagmcend%
  \tagpdfparaOn%
  \tagstructend%
  \vskip 0.25em%
  \tagmcbegin{artifact}%
    \hrule height 1pt%
  \tagmcend%
  \vskip 1em%
}
\makeatother
\title{Numerical Analysis PhD Exam, May 2020}
\author{}
\date{}
\begin{document}
\maketitle
\thispagestyle{empty}
\begin{examinstructions}
Do 4 (four) of the first 5 (1-5) and 4 (four) of the last 5 problems (6-10).
\end{examinstructions}
\begin{examproblems}{0}{H2}
\def\ProblemHeadingText{Problem 1.}
\item
This problem has two parts.
\begin{examsubparts}
\item[\textnormal{(a)}]
Let \(A \in \mathbb{C}^{m \times n}\). Prove or give a counterexample: Every matrix norm that is induced by a vector norm satisfies the submultiplicative property \(\|AB\| \leq \|A\|\|B\|\). If you prove this, make sure to justify each nontrivial step.
\item[\textnormal{(b)}]
Let \(A \in \mathbb{C}^{m \times m}\). Prove that \(\|A\|_2 = (\rho(A^*A))^{1/2}\), where \(\rho(A)\) is the spectral radius of~\(A\).
\end{examsubparts}
\def\ProblemHeadingText{Problem 2.}
\item
Suppose~\(A\) is Hermitian positive definite.
\begin{examsubparts}
\item[\textnormal{(a)}]
Prove that each principal submatrix of~\(A\) is Hermitian positive definite.
\item[\textnormal{(b)}]
Prove that an element of~\(A\) with largest magnitude lies on the diagonal.
\item[\textnormal{(c)}]
Prove that~\(A\) has a Cholesky decomposition.
\end{examsubparts}
\def\ProblemHeadingText{Problem 3.}
\item
Let \(A \in \mathbb{C}^{m \times m}\) be Hermitian.
\begin{examsubparts}
\item[\textnormal{(a)}]
Show that all eigenvalues of~\(A\) are real.
\item[\textnormal{(b)}]
Define the stationary iterative method (a.k.a. fixed point method)
\[
x^{(k+1)}=Ax^{(k)}+b. \tag{1}
\]
 Suppose (1) has fixed-point~\(x\), namely~\(x\) satisfies \(x=Ax+b\). Show the iteration (1) converges to~\(x\) from any starting guess \(x^{(0)}\), that is \(x^{(k)} \to x\) as \(k \to \infty\), if and only if the eigenvalues \(\lambda_i\) of~\(A\) satisfy \(|\lambda_i|<1\), \(i=1,\ldots,m\). You may use the fact that Hermitian matrix~\(A\) is unitarily diagonalizable.
\end{examsubparts}
\def\ProblemHeadingText{Problem 4.}
\item
Suppose the \(5 \times 5\) symmetric matrix~\(A\) has eigenvalues known to within the given tolerances.
\[
\begin{aligned}
3.5 &> \lambda_1 > 2.5\\
2.0 &> \lambda_2 > 1.0\\
1.0 &> \lambda_3 > -1.0\\
-1.0 &> \lambda_4 > -2.0\\
-2.5 &> \lambda_5 > -3.5.
\end{aligned}
\]
\begin{examsubparts}
\item[\textnormal{(a)}]
Describe how shifting can be used so that the power method can be used to compute \(\lambda_1\) with guaranteed convergence. Clearly explain your choice of shift.
\item[\textnormal{(b)}]
Provide an upper bound for the convergence rate using the shift you chose in (a) for \(\lambda_1\). Is there another shift that would decrease this worst-case convergence rate?
\end{examsubparts}
\def\ProblemHeadingText{Problem 5.}
\item
For \(x,y>0\), consider computing \(f(x,y)=\sqrt{y+x^2}-\sqrt{y}\) in floating-point arithmetic with machine precision \(\epsilon_m\).
\begin{examsubparts}
\item[\textnormal{(a)}]
Explain the difficulties in computing \(f(x,y)\), if \(x \ll y\). What are the absolute and relative errors if \(x^2/y<\epsilon_m\), if \(f(x,y)\) is computed directly from the form given above?
\item[\textnormal{(b)}]
Suppose \(x^2/y<\epsilon_m\). Describe a way to compute \(f(x,y)\) with more accuracy in this situation.
\end{examsubparts}
\def\ProblemHeadingText{Problem 6.}
\item
Let \(\alpha>0\). For (i) \(p=2\); (ii) \(p=\infty\), find the constant \(c_p\) that minimizes
\[
E_p(c)=\|t^\alpha-c\|_p=\left(\int_0^1 |t^\alpha-c|^p\,dt\right)^{1/p},
\]
 and find \(E_p(c_p)\), for each of those values of~\(p\).
\def\ProblemHeadingText{Problem 7.}
\item
Let \(\{\phi_k\}_{k=0}^{m+1}\) be the set of monic orthogonal polynomials with respect to inner product \((u,v)=\int_a^b u(x)v(x)w(x)\,dx\). Let \(\phi_{-1}=0\) and \(\phi_0=1\).
\begin{examsubparts}
\item[\textnormal{(a)}]
Find expressions for constants \(\alpha_k\) and \(\beta_k\) to determine a 3-term recurrence relation of the form
\[
\phi_{k+1}(x)=(x-\alpha_k)\phi_k-\beta_k\phi_{k-1}(x),\quad k=0,\ldots,m.
\]
\item[\textnormal{(b)}]
Use the above relation to determine \(\phi_1,\phi_2\) and \(\phi_3\), using the inner-product \((u,v)=\int_0^1 u(x)v(x)\,dx\).
\end{examsubparts}
\def\ProblemHeadingText{Problem 8.}
\item
Consider finding a zero of function \(f:D\subseteq\mathbb{R}^n\to\mathbb{R}^n\) that can be written as the sum of a linear and nonlinear part
\[
f(x)=Bx+G(x),
\]
 where~\(B\) is a nonsingular matrix and \(G:D\subseteq\mathbb{R}^n\to\mathbb{R}^n\) is some nonlinear function. At a point \(x_k\) consider an affine model \(M_k(x)=a_k+A_k(x-x_k)\), where the quantities \(a_k\in\mathbb{R}^n\) and \(A_k\in\mathbb{R}^{n\times n}\) are to be determined.
\begin{examsubparts}
\item[\textnormal{(a)}]
Determine \(a_k\) and \(A_k\) so that the following conditions hold:
\[
M_k(x_k)=f(x_k)\quad\text{and}\quad M_k'(x_k)=B.
\]
\item[\textnormal{(b)}]
Derive the iteration obtained by defining \(x_{k+1}\) as the zero of \(M_k(x)\).
\end{examsubparts}
\def\ProblemHeadingText{Problem 9.}
\item
Let \(f\in C^\infty(a,b)\). Let \(x_0<x_1<x_2\) be three points in \([a,b]\) that are not necessarily equally spaced.
\begin{examsubparts}
\item[\textnormal{(a)}]
Based on the quadratic interpolant \(p_2\) which satisfies \(p_2(x_0)=f(x_0)\), \(p_2(x_1)=f(x_1)\) and \(p_2(x_2)=f(x_2)\), find the centered finite difference approximations to \(f'(x_1)\) and \(f''(x_1)\) (you should explicitly show how the difference approximations are derived from the interpolant).
\item[\textnormal{(b)}]
Derive an expression for the error, \(f'(x_1)-p_2'(x_1)\).
\end{examsubparts}
\def\ProblemHeadingText{Problem 10.}
\item
This problem has two parts.
\begin{examsubparts}
\item[\textnormal{(a)}]
For \(f\in C^\infty(a,b)\), the composite trapezoidal quadrature rule with~\(n\) subintervals with length \(h=(b-a)/n\) satisfies
\[
\left|\int_a^b f(x)\,dx-I_{T,n}\right|=a_2h^2+a_4h^4+a_6h^6+\cdots,
\]
 where the coefficients \(a_2,a_4,\ldots\), do not depend on~\(n\). Find and inductively prove an expression for the~\(k\)th Richardson extrapolant \(I_k\) of \(I_0:=I_{T,n}\).
\item[\textnormal{(b)}]
Consider a quadrature formula of the type
\[
\int_0^1 f(x)\,dx\approx \alpha f(x_1)+\beta[f(1)-f(0)].
\]
 Determine \(\alpha,\beta\) and \(x_1\) such that the degree of exactness is as large as possible. What is the maximum degree of exactness?
\end{examsubparts}
\end{examproblems}
\end{document}
